% Copyright 2026 Nobin Sijo (NobinSijo7T).
% SPDX-License-Identifier: Apache-2.0
\documentclass[11pt,a4paper]{article}
\usepackage[margin=0.70in,top=0.68in,bottom=0.68in]{geometry}
\usepackage{amsmath,amssymb}
\usepackage{times}
\usepackage{enumitem}
\usepackage{fancyhdr}

\pagestyle{fancy}
\fancyhf{}
\fancyhead[L]{\small Group Project $\cdot$ PrismSpace Architecture}
\fancyhead[R]{\small\thepage}
\renewcommand{\headrulewidth}{0.4pt}
\setlength{\headheight}{14.5pt}

\setcounter{page}{6}

\setlength{\parindent}{0pt}
\setlength{\parskip}{0.28em}

\begin{document}

% =========================================================================
% PAGE 1 (Page 6): Overview, Inputs, Outputs, and Step 1
% =========================================================================

\subsection*{4.5\quad Algorithm of PrismSpace Multi-Agent System}

The PrismSpace algorithm orchestrates multiple specialized AI agents into a coordinated swarm: interpreting user objectives, routing tasks to the optimal LLM provider, dispatching specialized agents, equipping them with tools, and verifying safety and accuracy. The workflow operates in three distinct phases:
\begin{itemize}[leftmargin=1.5em, itemsep=0.10em, topsep=0.04em]
    \item \textbf{Phase I: Offline Swarm Graph Representation} (Steps 1--5: Data prep, swarm graph, GAT routing)
    \item \textbf{Phase II: Single-Stage Preference Alignment} (Step 6: Safety and quality tuning via ORPO)
    \item \textbf{Phase III: Online Runtime Inference \& Swarm Execution} (Steps 7--9: Live execution, evaluation, telemetry)
\end{itemize}

\vspace{0.20em}
\textbf{Inputs (What Goes In):}
\begin{itemize}[leftmargin=1.5em, itemsep=0.10em, topsep=0.04em]
    \item \textbf{User Objective ($\mathcal{T}_{\text{user}}$):} Natural language user prompt, operational scope, and execution constraints.
    \item \textbf{Curated Training Corpora ($\mathcal{D}$):} Multi-source benchmark datasets containing prompt traces, intent categories, safety evaluations (WildGuard), and multi-agent trajectories (RouterBench, LMSYS Arena, CMU Agent Benchmark).
    \item \textbf{Preference Alignment Pairs ($\mathcal{D}_{\text{pref}} = \{(x, y_w, y_l)\}$):} Prompt $x$ paired with preferred response $y_w$ and disfavored response $y_l$.
    \item \textbf{Global Agent Swarm Registry ($\mathcal{G}_{\text{global}} = (V, E)$):} Roster of specialized AI agents (Supervisor, Coding, Research, Validator, Planner) and Model Context Protocol (MCP) tool endpoints.
\end{itemize}

\vspace{0.20em}
\textbf{Outputs (What Comes Out):}
\begin{itemize}[leftmargin=1.5em, itemsep=0.10em, topsep=0.04em]
    \item \textbf{Task Intent \& Best Provider:} Predicted task category ($\hat{y}_{\text{intent}}$) and optimal LLM provider ($\hat{y}_{\text{provider}}$, e.g., Groq for rapid response, OpenAI/Anthropic for complex reasoning).
    \item \textbf{Swarm Dispatch Vector ($\hat{\mathbf{y}}_{\text{agents}}$):} Binary activation checklist ($1$ = active, $0$ = idle) for each agent in $V$.
    \item \textbf{Execution Plan ($\mathcal{G}_{\text{active}}$):} Directed Acyclic Graph (DAG) outlining task steps with predicted latency ($\hat{\ell}$) and cost ($\hat{c}$).
    \item \textbf{Tuned Model Adapters ($\Delta \mathbf{W}$):} Lightweight LoRA weights ensuring helpfulness and safety compliance.
    \item \textbf{Evaluation Metrics:} Quantitative benchmarks including Accuracy, F1-Score, ROC-AUC, and ECE calibration.
\end{itemize}

\vspace{0.35em}
\textbf{Step 1: Data collection and Preprocessing}
\begin{enumerate}[leftmargin=1.8em, itemsep=0.22em, topsep=0.06em]
    \item Import and unify raw records across multi-source benchmark datasets ($\mathcal{D}$), mapping heterogeneous prompt schemas, classification labels, and tool invocation logs into a structured, unified schema.
    \item Clean and sanitize the textual dataset by stripping malformed Unicode sequences, removing redundant whitespace tokens, trimming uninformative control characters, and filtering duplicate prompt-response pairs.
    \item Perform categorical encoding across discrete metadata attributes (agent role definitions, tool protocol capabilities, provider endpoints) into dense numerical vectors.
    \item Standardize continuous numerical features (token length, baseline query complexity, empirical latency traces) using $z$-score normalization: subtract the empirical feature mean and divide by standard deviation ($z = (x - \mu)/\sigma$) to yield zero mean and unit variance across all dimensions.
    \item Partition the processed dataset via stratified random sampling into three non-overlapping subsets:
    \begin{itemize}[leftmargin=1.5em, itemsep=0.08em, topsep=0.04em, label=$\ast$]
        \item 80\% Training set ($\mathcal{D}_{\text{train}}$) for model parameter optimization and graph embedding learning
        \item 10\% Validation set ($\mathcal{D}_{\text{val}}$) for hyperparameter tuning, pruning, and early stopping
        \item 10\% Testing set ($\mathcal{D}_{\text{test}}$) for unbiased evaluation of routing accuracy and safety alignment
    \end{itemize}
\end{enumerate}

\newpage

% =========================================================================
% PAGE 2 (Page 7): Steps 2, 3, 4, and 5
% =========================================================================

\textbf{Step 2: Swarm Graph Construction}
\begin{enumerate}[leftmargin=1.8em, itemsep=0.20em, topsep=0.05em]
    \item Define the node set ($V$): represent each specialized AI agent persona (Supervisor, Research, Coding, Validator, Planner) and each Model Context Protocol (MCP) tool as an individual vertex $v_i \in V$.
    \item Evaluate pairwise functional capability similarity: compute the Pearson correlation coefficient between node capability vectors $\mathbf{x}_i$ and $\mathbf{x}_j$, defined as the covariance of their feature distributions divided by the product of their standard deviations ($s_{ij} = \mathrm{Cov}(\mathbf{x}_i, \mathbf{x}_j) / (\sigma_i \sigma_j)$).
    \item Establish collaborative topological connections: generate binary adjacency matrix $\mathbf{A}$ by comparing similarity against correlation cutoff $\tau$, setting $A_{ij} = 1$ when pairwise similarity meets or exceeds threshold ($s_{ij} \ge \tau$) and $A_{ij} = 0$ otherwise.
    \item Construct the global swarm topology $\mathcal{G} = (V, E)$: establish undirected collaborative communication edges $E$ that permit message propagation exclusively between complementary agent and tool pairings.
\end{enumerate}

\vspace{0.30em}
\textbf{Step 3: Graph Attention Network (GAT) Initialization}
\begin{enumerate}[leftmargin=1.8em, itemsep=0.20em, topsep=0.05em]
    \item Initialize node representations: assign each vertex $v_i \in V$ an initial $d$-dimensional feature embedding vector $\mathbf{h}_i = \mathbf{x}_i \in \mathbb{R}^d$ describing agent operational scope, prompt semantics, and tool interface signatures.
    \item Initialize projection and attention parameters: instantiate weight projection matrix $\mathbf{W} \in \mathbb{R}^{d' \times d}$ and shared attention mechanism vector $\mathbf{a} \in \mathbb{R}^{2d'}$ via Xavier uniform initialization, ensuring stable variance across forward and backward signal propagation.
    \item Configure network hyperparameters: establish learning rate $\eta$, attention head count $K$ (specialized across task intent, model capability, swarm affinity, and safety constraints), and dropout regularization probability $p$.
\end{enumerate}

\vspace{0.30em}
\textbf{Step 4: Attention-Based Message Passing}

For each node $v_i$ with immediate collaborative neighborhood $\mathcal{N}(i)$:
\begin{enumerate}[leftmargin=1.8em, itemsep=0.20em, topsep=0.05em]
    \item Compute raw attention coefficients with connected neighbors $v_j \in \mathcal{N}(i)$: project node features through transformation matrix $\mathbf{W}$, concatenate projected states $[\mathbf{W}\mathbf{h}_i \parallel \mathbf{W}\mathbf{h}_j]$, compute inner product with attention vector $\mathbf{a}$, and apply LeakyReLU non-linearity: $e_{ij} = \mathrm{LeakyReLU}(\mathbf{a}^T [\mathbf{W} \mathbf{h}_i \parallel \mathbf{W} \mathbf{h}_j])$.
    \item Normalize coefficients across the local neighborhood using Softmax: calculate positive attention weights $\alpha_{ij} = \exp(e_{ij}) / \sum_{k \in \mathcal{N}(i)} \exp(e_{ik})$ that sum to unity, quantifying the relative communicative significance of neighbor $j$ to node $i$.
    \item Aggregate neighborhood feature representations: compute a weighted linear sum of neighbors' transformed representations scaled by attention weights, followed by non-linear activation $\mathbf{h}_i' = \sigma\big(\sum_{j \in \mathcal{N}(i)} \alpha_{ij} \mathbf{W} \mathbf{h}_j\big)$.
    \item Multi-head attention concatenation: execute $K$ independent attention heads in parallel and concatenate their output representations $\mathbf{h}_i' = \Vert_{k=1}^K \sigma\big(\sum_{j \in \mathcal{N}(i)} \alpha_{ij}^k \mathbf{W}^k \mathbf{h}_j\big)$, capturing rich, multi-perspective collaborative semantics.
\end{enumerate}

\vspace{0.30em}
\textbf{Step 5: Classification and Loss Computation}
\begin{enumerate}[leftmargin=1.8em, itemsep=0.20em, topsep=0.05em]
    \item Route updated swarm representations $\mathbf{h}_i'$ through dual specialized prediction heads:
    \begin{itemize}[leftmargin=1.5em, itemsep=0.06em, topsep=0.03em, label=$\ast$]
        \item \textit{Intent \& Provider Routing Head:} Computes normalized routing probabilities via Softmax ($\hat{\mathbf{y}}_{\text{route}} = \mathrm{Softmax}(\mathbf{W}_{\text{route}} \mathbf{h}' + \mathbf{b}_{\text{route}})$), selecting the optimal foundation model provider based on task difficulty and speed.
        \item \textit{Agent Dispatch Head:} Evaluates independent activation probabilities per agent persona via element-wise Sigmoid ($\hat{y}_{\text{agent}, i} = \mathrm{Sigmoid}(\mathbf{W}_{\text{agent}} \mathbf{h}_i' + b_i)$), determining active swarm membership.
    \end{itemize}
    \item Compute unified multi-task loss ($\mathcal{L}_{\text{class}}$): combine categorical cross-entropy for provider routing with binary cross-entropy across agent dispatch activations, balancing classification accuracy with minimal swarm over-provisioning: $\mathcal{L}_{\text{class}} = \mathcal{L}_{\text{CE}}(\mathbf{y}_{\text{route}}, \hat{\mathbf{y}}_{\text{route}}) + \lambda \sum_{v \in V} \mathcal{L}_{\text{BCE}}(y_v, \hat{y}_v)$.
    \item Optimize parameters: backpropagate error gradients and update network parameters using the AdamW optimizer with cosine learning-rate decay and gradient norm clipping.
\end{enumerate}

\newpage

% =========================================================================
% PAGE 3 (Page 8): Steps 6 and 7
% =========================================================================

\textbf{Step 6: Single-Stage Preference Alignment (ORPO)}
\begin{enumerate}[leftmargin=1.8em, itemsep=0.22em, topsep=0.06em]
    \item Inject Parameter-Efficient Fine-Tuning (PEFT) adapters: attach low-rank decomposition matrices ($\Delta \mathbf{W} = \frac{\alpha}{r} \mathbf{B}\mathbf{A}$) to attention projection layers of the base language model, enabling targeted adaptation while keeping foundational weights frozen.
    \item Compute the supervised fine-tuning (SFT) objective: calculate the sequence negative log-likelihood loss exclusively on favored human responses ($y_w$), preserving foundational language fluency and structured syntax: $\mathcal{L}_{\mathrm{SFT}} = -\frac{1}{|y_w|} \sum_{t=1}^{|y_w|} \log P(y_{w,t} \mid x, y_{w,<t})$.
    \item Calculate length-normalized sequence odds: evaluate the likelihood of token sequence $y \in \{y_w, y_l\}$ using the odds formulation $\mathrm{odds}(y \mid x) = P(y \mid x) / [1 - P(y \mid x)]$, where per-token probabilities are geometrically averaged across sequence length to eliminate bias favoring shorter answers.
    \item Formulate the Odds Ratio (OR) preference penalty: calculate the negative log-sigmoid of the log-odds ratio between favored ($y_w$) and disfavored ($y_l$) completions: $\mathcal{L}_{\mathrm{OR}} = -\log \sigma \big(\log [\mathrm{odds}(y_w \mid x) / \mathrm{odds}(y_l \mid x)]\big)$, actively disincentivizing unsafe, hallucinated, or low-quality patterns.
    \item Optimize unified single-stage ORPO objective: jointly minimize the composite loss $\mathcal{L}_{\mathrm{ORPO}} = \mathcal{L}_{\mathrm{SFT}} + \beta \cdot \mathcal{L}_{\mathrm{OR}}$, achieving concurrent instruction following and safety preference alignment without requiring a separate reward model or reinforcement learning phase.
\end{enumerate}

\vspace{0.40em}
\textbf{Step 7: Real-Time Runtime Inference and Swarm Execution}
\begin{enumerate}[leftmargin=1.8em, itemsep=0.22em, topsep=0.06em]
    \item Capture user objective: ingest natural language prompt $\mathcal{T}_{\text{user}}$ via the Next.js client interface and securely dispatch it to the FastAPI Hive Bridge orchestration gateway.
    \item Synthesize query feature vector ($\mathbf{x}_{\text{query}}$): concatenate sparse TF-IDF n-gram vectors, dense semantic sentence embeddings, and session contextual metadata (token budget, latency requirements, user tool permissions).
    \item Determine routing and dispatch policy: assign the prompt to the highest-scoring LLM provider via argmax selection on routing probabilities ($\hat{y}_{\text{provider}} = \mathrm{argmax}(\hat{\mathbf{y}}_{\text{route}})$), and activate all worker agents whose dispatch score meets threshold $\tau_{\text{dispatch}}$ ($\hat{\mathbf{y}}_{\text{agents}} = \mathbb{I}(\hat{\mathbf{y}}_{\text{dispatch}} \ge \tau_{\text{dispatch}})$).
    \item Formulate task DAG and execution schedule: generate active dependency graph $\mathcal{G}_{\text{active}}$ and execute a topological sort ($\pi = \mathrm{TopologicalSort}(\mathcal{G}_{\text{active}})$) to determine parallel and sequential task execution phases.
    \item Execute autonomous ReAct swarm loops: dispatch agents to iteratively reason, propose actions, and invoke Model Context Protocol (MCP) tool endpoints (Browser, Filesystem, GitHub, Python Sandbox), aggregating intermediate tool observation traces until all subtask criteria are satisfied.
    \item Enforce policy safety gating: evaluate final candidate completions through the calibrated safety gate, verifying that the log-odds safety score relative to reference text meets safety threshold $\tau_{\text{safety}}$, intercepting prompt injection, unsafe operations, or policy violations.
    \item Deliver real-time response and persist session state: stream verified response tokens live to the frontend via Server-Sent Events (SSE), update real-time telemetry HUD metrics, and commit complete conversation artifacts to local browser IndexedDB storage using Dexie.js.
\end{enumerate}

\newpage

% =========================================================================
% PAGE 4 (Page 9): Steps 8, 9, and Summary
% =========================================================================

\textbf{Step 8: Model Evaluation}
\begin{enumerate}[leftmargin=1.8em, itemsep=0.22em, topsep=0.06em]
    \item Assess generalization performance: benchmark trained GAT routing and ORPO alignment models against holdout validation ($\mathcal{D}_{\text{val}}$) and test ($\mathcal{D}_{\text{test}}$) splits to quantify accuracy and preference adherence.
    \item Compute comprehensive quantitative validation metrics:
    \begin{itemize}[leftmargin=1.5em, itemsep=0.08em, topsep=0.04em, label=$\ast$]
        \item \textit{Routing Accuracy \& Macro-F1:} Measure proportion of correctly classified intents and provider selections, computing harmonic mean of precision and recall to evaluate balanced performance across rare and frequent task classes.
        \item \textit{Discrimination Capability (ROC-AUC \& AUPRC):} Measure the area under the receiver operating characteristic and precision--recall curves to assess safety gate discrimination under class imbalance.
        \item \textit{Calibration Quality (ECE):} Measure Expected Calibration Error across $M$ probability bins ($\mathrm{ECE} = \sum_{m=1}^M \frac{|B_m|}{N} |\mathrm{acc}(B_m) - \mathrm{conf}(B_m)|$) to verify that predicted routing confidence aligns with empirical accuracy.
        \item \textit{Regression Precision ($R^2$ \& MAE):} Quantify predictive fidelity of task latency ($\hat{\ell}$) and monetary token cost ($\hat{c}$) against actual multi-agent execution traces.
    \end{itemize}
    \item Save optimized model checkpoints, export weight artifacts, and log diagnostic performance reports.
\end{enumerate}

\vspace{0.40em}
\textbf{Step 9: Visualization and Interpretation}
\begin{itemize}[leftmargin=1.8em, itemsep=0.22em, topsep=0.06em, label=$\blacktriangleright$]
    \item Generate convergence diagnostics: plot training and validation loss trajectories, ROC curves, and precision--recall frontiers to track learning progression and over-fitting boundaries.
    \item Inspect attention heatmaps: visualize inter-agent and agent-tool attention weight matrices across the $K$ attention heads, identifying dominant collaboration pathways and tool utilization dependencies.
    \item Monitor live operational telemetry: track real-time routing hit rates, token cost allocations ($\hat{c}$), task latency profiles ($\hat{\ell}$), and safety policy pass rates on the Swarm Telemetry HUD.
\end{itemize}

\vspace{0.40em}
\textbf{Algorithm Summary}
\begin{enumerate}[leftmargin=1.8em, itemsep=0.18em, topsep=0.05em]
    \item \textbf{Graph-Attentive Swarm Routing:} The GAT architecture models functional interactions between specialized agents and MCP tools, dynamically routing incoming tasks to the optimal LLM provider and dispatching minimal required worker agents.
    \item \textbf{Single-Stage Preference Alignment:} The ORPO formulation unifies supervised fine-tuning with length-normalized odds ratio penalties, tuning model helpfulness and safety in a single stage without independent reward models.
    \item \textbf{Autonomous Reactive Execution:} Runtime execution is scheduled via DAG topological sorting, iterative ReAct tool loops, rigorous odds-ratio safety gating, and live SSE streaming to browser IndexedDB storage.
\end{enumerate}

\end{document}
